% Options for packages loaded elsewhere
% Options for packages loaded elsewhere
\PassOptionsToPackage{unicode}{hyperref}
\PassOptionsToPackage{hyphens}{url}
\PassOptionsToPackage{dvipsnames,svgnames,x11names}{xcolor}
%
\documentclass[
  letterpaper,
  DIV=11,
  numbers=noendperiod]{scrartcl}
\usepackage{xcolor}
\usepackage{amsmath,amssymb}
\setcounter{secnumdepth}{5}
\usepackage{iftex}
\ifPDFTeX
  \usepackage[T1]{fontenc}
  \usepackage[utf8]{inputenc}
  \usepackage{textcomp} % provide euro and other symbols
\else % if luatex or xetex
  \usepackage{unicode-math} % this also loads fontspec
  \defaultfontfeatures{Scale=MatchLowercase}
  \defaultfontfeatures[\rmfamily]{Ligatures=TeX,Scale=1}
\fi
\usepackage{lmodern}
\ifPDFTeX\else
  % xetex/luatex font selection
\fi
% Use upquote if available, for straight quotes in verbatim environments
\IfFileExists{upquote.sty}{\usepackage{upquote}}{}
\IfFileExists{microtype.sty}{% use microtype if available
  \usepackage[]{microtype}
  \UseMicrotypeSet[protrusion]{basicmath} % disable protrusion for tt fonts
}{}
\makeatletter
\@ifundefined{KOMAClassName}{% if non-KOMA class
  \IfFileExists{parskip.sty}{%
    \usepackage{parskip}
  }{% else
    \setlength{\parindent}{0pt}
    \setlength{\parskip}{6pt plus 2pt minus 1pt}}
}{% if KOMA class
  \KOMAoptions{parskip=half}}
\makeatother
% Make \paragraph and \subparagraph free-standing
\makeatletter
\ifx\paragraph\undefined\else
  \let\oldparagraph\paragraph
  \renewcommand{\paragraph}{
    \@ifstar
      \xxxParagraphStar
      \xxxParagraphNoStar
  }
  \newcommand{\xxxParagraphStar}[1]{\oldparagraph*{#1}\mbox{}}
  \newcommand{\xxxParagraphNoStar}[1]{\oldparagraph{#1}\mbox{}}
\fi
\ifx\subparagraph\undefined\else
  \let\oldsubparagraph\subparagraph
  \renewcommand{\subparagraph}{
    \@ifstar
      \xxxSubParagraphStar
      \xxxSubParagraphNoStar
  }
  \newcommand{\xxxSubParagraphStar}[1]{\oldsubparagraph*{#1}\mbox{}}
  \newcommand{\xxxSubParagraphNoStar}[1]{\oldsubparagraph{#1}\mbox{}}
\fi
\makeatother

\usepackage{color}
\usepackage{fancyvrb}
\newcommand{\VerbBar}{|}
\newcommand{\VERB}{\Verb[commandchars=\\\{\}]}
\DefineVerbatimEnvironment{Highlighting}{Verbatim}{commandchars=\\\{\}}
% Add ',fontsize=\small' for more characters per line
\usepackage{framed}
\definecolor{shadecolor}{RGB}{241,243,245}
\newenvironment{Shaded}{\begin{snugshade}}{\end{snugshade}}
\newcommand{\AlertTok}[1]{\textcolor[rgb]{0.68,0.00,0.00}{#1}}
\newcommand{\AnnotationTok}[1]{\textcolor[rgb]{0.37,0.37,0.37}{#1}}
\newcommand{\AttributeTok}[1]{\textcolor[rgb]{0.40,0.45,0.13}{#1}}
\newcommand{\BaseNTok}[1]{\textcolor[rgb]{0.68,0.00,0.00}{#1}}
\newcommand{\BuiltInTok}[1]{\textcolor[rgb]{0.00,0.23,0.31}{#1}}
\newcommand{\CharTok}[1]{\textcolor[rgb]{0.13,0.47,0.30}{#1}}
\newcommand{\CommentTok}[1]{\textcolor[rgb]{0.37,0.37,0.37}{#1}}
\newcommand{\CommentVarTok}[1]{\textcolor[rgb]{0.37,0.37,0.37}{\textit{#1}}}
\newcommand{\ConstantTok}[1]{\textcolor[rgb]{0.56,0.35,0.01}{#1}}
\newcommand{\ControlFlowTok}[1]{\textcolor[rgb]{0.00,0.23,0.31}{\textbf{#1}}}
\newcommand{\DataTypeTok}[1]{\textcolor[rgb]{0.68,0.00,0.00}{#1}}
\newcommand{\DecValTok}[1]{\textcolor[rgb]{0.68,0.00,0.00}{#1}}
\newcommand{\DocumentationTok}[1]{\textcolor[rgb]{0.37,0.37,0.37}{\textit{#1}}}
\newcommand{\ErrorTok}[1]{\textcolor[rgb]{0.68,0.00,0.00}{#1}}
\newcommand{\ExtensionTok}[1]{\textcolor[rgb]{0.00,0.23,0.31}{#1}}
\newcommand{\FloatTok}[1]{\textcolor[rgb]{0.68,0.00,0.00}{#1}}
\newcommand{\FunctionTok}[1]{\textcolor[rgb]{0.28,0.35,0.67}{#1}}
\newcommand{\ImportTok}[1]{\textcolor[rgb]{0.00,0.46,0.62}{#1}}
\newcommand{\InformationTok}[1]{\textcolor[rgb]{0.37,0.37,0.37}{#1}}
\newcommand{\KeywordTok}[1]{\textcolor[rgb]{0.00,0.23,0.31}{\textbf{#1}}}
\newcommand{\NormalTok}[1]{\textcolor[rgb]{0.00,0.23,0.31}{#1}}
\newcommand{\OperatorTok}[1]{\textcolor[rgb]{0.37,0.37,0.37}{#1}}
\newcommand{\OtherTok}[1]{\textcolor[rgb]{0.00,0.23,0.31}{#1}}
\newcommand{\PreprocessorTok}[1]{\textcolor[rgb]{0.68,0.00,0.00}{#1}}
\newcommand{\RegionMarkerTok}[1]{\textcolor[rgb]{0.00,0.23,0.31}{#1}}
\newcommand{\SpecialCharTok}[1]{\textcolor[rgb]{0.37,0.37,0.37}{#1}}
\newcommand{\SpecialStringTok}[1]{\textcolor[rgb]{0.13,0.47,0.30}{#1}}
\newcommand{\StringTok}[1]{\textcolor[rgb]{0.13,0.47,0.30}{#1}}
\newcommand{\VariableTok}[1]{\textcolor[rgb]{0.07,0.07,0.07}{#1}}
\newcommand{\VerbatimStringTok}[1]{\textcolor[rgb]{0.13,0.47,0.30}{#1}}
\newcommand{\WarningTok}[1]{\textcolor[rgb]{0.37,0.37,0.37}{\textit{#1}}}

\usepackage{longtable,booktabs,array}
\usepackage{calc} % for calculating minipage widths
% Correct order of tables after \paragraph or \subparagraph
\usepackage{etoolbox}
\makeatletter
\patchcmd\longtable{\par}{\if@noskipsec\mbox{}\fi\par}{}{}
\makeatother
% Allow footnotes in longtable head/foot
\IfFileExists{footnotehyper.sty}{\usepackage{footnotehyper}}{\usepackage{footnote}}
\makesavenoteenv{longtable}
\usepackage{graphicx}
\makeatletter
\newsavebox\pandoc@box
\newcommand*\pandocbounded[1]{% scales image to fit in text height/width
  \sbox\pandoc@box{#1}%
  \Gscale@div\@tempa{\textheight}{\dimexpr\ht\pandoc@box+\dp\pandoc@box\relax}%
  \Gscale@div\@tempb{\linewidth}{\wd\pandoc@box}%
  \ifdim\@tempb\p@<\@tempa\p@\let\@tempa\@tempb\fi% select the smaller of both
  \ifdim\@tempa\p@<\p@\scalebox{\@tempa}{\usebox\pandoc@box}%
  \else\usebox{\pandoc@box}%
  \fi%
}
% Set default figure placement to htbp
\def\fps@figure{htbp}
\makeatother





\setlength{\emergencystretch}{3em} % prevent overfull lines

\providecommand{\tightlist}{%
  \setlength{\itemsep}{0pt}\setlength{\parskip}{0pt}}



 


% load packages
\usepackage{geometry}
\usepackage{xcolor}
\usepackage{eso-pic}
\usepackage{fancyhdr}
\usepackage{sectsty}
\usepackage{fontspec}
\usepackage{titlesec}

%% Set page size with a wider right margin
\geometry{a4paper, total={170mm,257mm}, left=20mm, top=20mm, bottom=20mm, right=50mm}

%% Let's define some colours
\definecolor{light}{HTML}{E6E6FA}
\definecolor{highlight}{HTML}{800080}
\definecolor{dark}{HTML}{330033}

%% Let's add the border on the right hand side 
\AddToShipoutPicture{% 
    \AtPageLowerLeft{% 
        \put(\LenToUnit{\dimexpr\paperwidth-3cm},0){% 
            \color{light}\rule{3cm}{\LenToUnit\paperheight}%
          }%
     }%
     % logo
    \AtPageLowerLeft{% start the bar at the bottom right of the page
        \put(\LenToUnit{\dimexpr\paperwidth-2.25cm},27.2cm){% move it to the top right
            \color{light}
\includegraphics[width=1.5cm]{_extensions/nrennie/PrettyPDF/logo.png}
          }%
     }%
}

%% Style the page number
\fancypagestyle{mystyle}{
  \fancyhf{}
  \renewcommand\headrulewidth{0pt}
  \fancyfoot[R]{\thepage}
  \fancyfootoffset{3.5cm}
}
\setlength{\footskip}{20pt}

%% style the chapter/section fonts
\chapterfont{\color{dark}\fontsize{20}{16.8}\selectfont}
\sectionfont{\color{dark}\fontsize{20}{16.8}\selectfont}
\subsectionfont{\color{dark}\fontsize{14}{16.8}\selectfont}
\titleformat{\subsection}
  {\sffamily\Large\bfseries}{\thesection}{1em}{}[{\titlerule[0.8pt]}]
  
% left align title
\makeatletter
\renewcommand{\maketitle}{\bgroup\setlength{\parindent}{0pt}
\begin{flushleft}
  {\sffamily\huge\textbf{\MakeUppercase{\@title}}} \vspace{0.3cm} \newline
  {\Large {\@subtitle}} \newline
  \@author
\end{flushleft}\egroup
}
\makeatother

%% Use some custom fonts
\setsansfont{Ubuntu}[
    Path=_extensions/nrennie/PrettyPDF/Ubuntu/,
    Scale=0.9,
    Extension = .ttf,
    UprightFont=*-Regular,
    BoldFont=*-Bold,
    ItalicFont=*-Italic,
    ]

\setmainfont{Ubuntu}[
    Path=_extensions/nrennie/PrettyPDF/Ubuntu/,
    Scale=0.9,
    Extension = .ttf,
    UprightFont=*-Regular,
    BoldFont=*-Bold,
    ItalicFont=*-Italic,
    ]
\KOMAoption{captions}{tableheading}
\makeatletter
\@ifpackageloaded{tcolorbox}{}{\usepackage[skins,breakable]{tcolorbox}}
\@ifpackageloaded{fontawesome5}{}{\usepackage{fontawesome5}}
\definecolor{quarto-callout-color}{HTML}{909090}
\definecolor{quarto-callout-note-color}{HTML}{0758E5}
\definecolor{quarto-callout-important-color}{HTML}{CC1914}
\definecolor{quarto-callout-warning-color}{HTML}{EB9113}
\definecolor{quarto-callout-tip-color}{HTML}{00A047}
\definecolor{quarto-callout-caution-color}{HTML}{FC5300}
\definecolor{quarto-callout-color-frame}{HTML}{acacac}
\definecolor{quarto-callout-note-color-frame}{HTML}{4582ec}
\definecolor{quarto-callout-important-color-frame}{HTML}{d9534f}
\definecolor{quarto-callout-warning-color-frame}{HTML}{f0ad4e}
\definecolor{quarto-callout-tip-color-frame}{HTML}{02b875}
\definecolor{quarto-callout-caution-color-frame}{HTML}{fd7e14}
\makeatother
\makeatletter
\@ifpackageloaded{caption}{}{\usepackage{caption}}
\AtBeginDocument{%
\ifdefined\contentsname
  \renewcommand*\contentsname{Table of contents}
\else
  \newcommand\contentsname{Table of contents}
\fi
\ifdefined\listfigurename
  \renewcommand*\listfigurename{List of Figures}
\else
  \newcommand\listfigurename{List of Figures}
\fi
\ifdefined\listtablename
  \renewcommand*\listtablename{List of Tables}
\else
  \newcommand\listtablename{List of Tables}
\fi
\ifdefined\figurename
  \renewcommand*\figurename{Figure}
\else
  \newcommand\figurename{Figure}
\fi
\ifdefined\tablename
  \renewcommand*\tablename{Table}
\else
  \newcommand\tablename{Table}
\fi
}
\@ifpackageloaded{float}{}{\usepackage{float}}
\floatstyle{ruled}
\@ifundefined{c@chapter}{\newfloat{codelisting}{h}{lop}}{\newfloat{codelisting}{h}{lop}[chapter]}
\floatname{codelisting}{Listing}
\newcommand*\listoflistings{\listof{codelisting}{List of Listings}}
\makeatother
\makeatletter
\makeatother
\makeatletter
\@ifpackageloaded{caption}{}{\usepackage{caption}}
\@ifpackageloaded{subcaption}{}{\usepackage{subcaption}}
\makeatother
\makeatletter
\@ifpackageloaded{tcolorbox}{}{\usepackage[skins,breakable]{tcolorbox}}
\makeatother
\makeatletter
\@ifundefined{shadecolor}{\definecolor{shadecolor}{rgb}{.97, .97, .97}}{}
\makeatother
\makeatletter
\@ifundefined{codebgcolor}{\definecolor{codebgcolor}{named}{light}}{}
\makeatother
\makeatletter
\ifdefined\Shaded\renewenvironment{Shaded}{\begin{tcolorbox}[frame hidden, boxrule=0pt, sharp corners, enhanced, breakable, colback={codebgcolor}]}{\end{tcolorbox}}\fi
\makeatother
\usepackage{bookmark}
\IfFileExists{xurl.sty}{\usepackage{xurl}}{} % add URL line breaks if available
\urlstyle{same}
\hypersetup{
  pdftitle={Practical 2 - GLMM},
  colorlinks=true,
  linkcolor={highlight},
  filecolor={Maroon},
  citecolor={Blue},
  urlcolor={highlight},
  pdfcreator={LaTeX via pandoc}}


\title{Practical 2 - GLMM}
\author{}
\date{}
\begin{document}
\maketitle

\pagestyle{mystyle}


In this practical we are going to fit a Generalized (Mixed) Linear Model
in \texttt{inlabru}.

We are going to:

\begin{itemize}
\tightlist
\item
  Fit a Poisson regression
\item
  Change the prior distributions for the model hyperparameters (both for
  fixed and random effects)
\item
  Compute and visualize posterior densities and summaries for marginal
  effects
\end{itemize}

We start by loading some useful libraries

\begin{Shaded}
\begin{Highlighting}[]
\FunctionTok{library}\NormalTok{(dplyr)}
\FunctionTok{library}\NormalTok{(INLA)}
\FunctionTok{library}\NormalTok{(ggplot2)}
\FunctionTok{library}\NormalTok{(patchwork)}
\FunctionTok{library}\NormalTok{(inlabru)     }
\FunctionTok{library}\NormalTok{(lme4)}
\end{Highlighting}
\end{Shaded}

In this practical we are going to analyse the dataset
\texttt{grouseticks}. The data, contained in the library \texttt{lme4},
Number of ticks on the heads of red grouse chicks sampled in the field.
See \texttt{?grouseticks} for details.

The data are analysed in the paper

Elston et al.~``Analysis of aggregation, a worked example: numbers of
ticks on red grouse chicks.'' Parasitology

and in this practical we will follow their analysis.

\begin{Shaded}
\begin{Highlighting}[]
\FunctionTok{data}\NormalTok{(grouseticks)}
\NormalTok{?grouseticks}
\FunctionTok{head}\NormalTok{(grouseticks)}
\end{Highlighting}
\end{Shaded}

\begin{verbatim}
  INDEX TICKS BROOD HEIGHT YEAR LOCATION   cHEIGHT
1     1     0   501    465   95       32  2.759305
2     2     0   501    465   95       32  2.759305
3     3     0   502    472   95       36  9.759305
4     4     0   503    475   95       37 12.759305
5     5     0   503    475   95       37 12.759305
6     6     3   503    475   95       37 12.759305
\end{verbatim}

\subsection{Fitting Poisson
regression}\label{fitting-poisson-regression}

We assume that the number of ticks \(y_{ijk}\) counted on chick \(i\) of
brood \(j\) in year \(k\) follows a Poisson distribution with mean
\(\lambda_{ijk}\) \[
y_{ijk}|\lambda_{ijk}\sim\text{Poisson}(\lambda_{ijk})
\] We then model log mean counts \(\eta_{ijk} = \log(\lambda_{ijk})\) as
a linear function of year, altitude, brood, and individual chick within
brood. We assume a fixed effect \(\alpha_k\) of year \(k\), a linear
effect of altitude \(x_{ij}\), two random effects \(e_{jk}\) and
\(\epsilon_{ijk}\) brood and individual within brood respectively. Thus:
\begin{equation}\phantomsection\label{eq-linpred}{
\eta_{ijk} = \log(\lambda_{ijk}) = \alpha_k + \beta x_{ij} + e_{jk} + \epsilon_{ijk}
}\end{equation}

where \(e_{jk}\sim\mathcal{N}(0,\tau^{-1}_e)\) and
\(\epsilon_{jk}\sim\mathcal{N}(0,\tau^{-1}_\epsilon)\).

We first fit the model using the default priors for fixed and random
effects.

To fit the model we first have to create two more variable, one that
indexes the combination \(jk\) of brood and year and another that
indexes the combination \(ijk\) of individuals per brood per year.

\begin{Shaded}
\begin{Highlighting}[]
\NormalTok{grouseticks  }\OtherTok{=}\NormalTok{ grouseticks }\SpecialCharTok{\%\textgreater{}\%}
  \FunctionTok{group\_by}\NormalTok{(BROOD, YEAR) }\SpecialCharTok{\%\textgreater{}\%}
  \FunctionTok{mutate}\NormalTok{(}\AttributeTok{ij =} \FunctionTok{cur\_group\_id}\NormalTok{()) }\SpecialCharTok{\%\textgreater{}\%}
  \FunctionTok{ungroup}\NormalTok{() }\SpecialCharTok{\%\textgreater{}\%}
  \FunctionTok{mutate}\NormalTok{(}\AttributeTok{ijk =} \FunctionTok{seq\_along}\NormalTok{(INDEX))}
\end{Highlighting}
\end{Shaded}

now we can fit the model.

\begin{tcolorbox}[enhanced jigsaw, opacityback=0, titlerule=0mm, toptitle=1mm, toprule=.15mm, opacitybacktitle=0.6, colframe=quarto-callout-warning-color-frame, arc=.35mm, rightrule=.15mm, title={Task}, colback=white, bottomrule=.15mm, bottomtitle=1mm, coltitle=black, left=2mm, leftrule=.75mm, breakable, colbacktitle=quarto-callout-warning-color!10!white]

Complete the code to fit the model in Equation~\ref{eq-linpred}. You
need to define the components, the likelihood and finaly use the
\texttt{bru()} function to get the results

\begin{Shaded}
\begin{Highlighting}[]
\NormalTok{cmp}\OtherTok{=} \ErrorTok{\textasciitilde{}} \SpecialCharTok{{-}}\DecValTok{1} \SpecialCharTok{+} \FunctionTok{year}\NormalTok{(...) }\SpecialCharTok{+}
  \FunctionTok{height}\NormalTok{(...) }\SpecialCharTok{+} 
  \FunctionTok{brood\_year}\NormalTok{(...) }\SpecialCharTok{+}
  \FunctionTok{brood\_year\_chicken}\NormalTok{(...)}

\NormalTok{lik }\OtherTok{=} \FunctionTok{bru\_obs}\NormalTok{(}\AttributeTok{forumula =}\NormalTok{ ...,}
              \AttributeTok{data =}\NormalTok{ ...,}
              \AttributeTok{family =}\NormalTok{ ...)}

\NormalTok{fit }\OtherTok{=} \FunctionTok{bru}\NormalTok{(cmp, lik)}
\end{Highlighting}
\end{Shaded}

Click here to see the solution

\begin{Shaded}
\begin{Highlighting}[]
\NormalTok{cmp}\OtherTok{=} \ErrorTok{\textasciitilde{}} \SpecialCharTok{{-}}\DecValTok{1} \SpecialCharTok{+} \FunctionTok{year}\NormalTok{(YEAR, }\AttributeTok{model =} \StringTok{"iid"}\NormalTok{, }\AttributeTok{initial =} \FunctionTok{log}\NormalTok{(}\FloatTok{0.01}\NormalTok{), }\AttributeTok{fixed =}\NormalTok{ T) }\SpecialCharTok{+}
  \FunctionTok{height}\NormalTok{(cHEIGHT, }\AttributeTok{model =}  \StringTok{"linear"}\NormalTok{) }\SpecialCharTok{+} 
  \FunctionTok{brood\_year}\NormalTok{(ij , }\AttributeTok{model =} \StringTok{"iid"}\NormalTok{) }\SpecialCharTok{+}
  \FunctionTok{brood\_year\_chicken}\NormalTok{(ijk, }\AttributeTok{model=} \StringTok{"iid"}\NormalTok{)}

\NormalTok{lik }\OtherTok{=} \FunctionTok{bru\_obs}\NormalTok{(TICKS }\SpecialCharTok{\textasciitilde{}}\NormalTok{ .,}
              \AttributeTok{data =}\NormalTok{ grouseticks,}
              \AttributeTok{family =} \StringTok{"Poisson"}\NormalTok{)}

  
\NormalTok{fit }\OtherTok{=} \FunctionTok{bru}\NormalTok{(cmp, lik)}
\end{Highlighting}
\end{Shaded}

\end{tcolorbox}

Now we can check the results.

\begin{tcolorbox}[enhanced jigsaw, opacityback=0, titlerule=0mm, toptitle=1mm, toprule=.15mm, opacitybacktitle=0.6, colframe=quarto-callout-warning-color-frame, arc=.35mm, rightrule=.15mm, title={Task}, colback=white, bottomrule=.15mm, bottomtitle=1mm, coltitle=black, left=2mm, leftrule=.75mm, breakable, colbacktitle=quarto-callout-warning-color!10!white]

Plot the posterior distribution of the fixed effects: \(\beta\) and
\(\alpha_k\).

Click here to see the solution

\begin{Shaded}
\begin{Highlighting}[]
\NormalTok{fit}\SpecialCharTok{$}\NormalTok{marginals.fixed}\SpecialCharTok{$}\NormalTok{height }\SpecialCharTok{\%\textgreater{}\%} \FunctionTok{ggplot}\NormalTok{() }\SpecialCharTok{+} \FunctionTok{geom\_line}\NormalTok{(}\FunctionTok{aes}\NormalTok{(x,y)) }\SpecialCharTok{+} \FunctionTok{ggtitle}\NormalTok{(}\StringTok{"Linear effect of altitude"}\NormalTok{)}
\end{Highlighting}
\end{Shaded}

\pandocbounded{\includegraphics[keepaspectratio]{GLMM_fit_files/figure-pdf/unnamed-chunk-7-1.pdf}}

\begin{Shaded}
\begin{Highlighting}[]
\FunctionTok{ggplot}\NormalTok{() }\SpecialCharTok{+} \FunctionTok{geom\_line}\NormalTok{(}\AttributeTok{data =}\NormalTok{ fit}\SpecialCharTok{$}\NormalTok{marginals.random}\SpecialCharTok{$}\NormalTok{year[[}\DecValTok{1}\NormalTok{]], }\FunctionTok{aes}\NormalTok{(x,y, }\AttributeTok{color =} \StringTok{"96"}\NormalTok{) )}\SpecialCharTok{+}
  \FunctionTok{geom\_line}\NormalTok{(}\AttributeTok{data =}\NormalTok{ fit}\SpecialCharTok{$}\NormalTok{marginals.random}\SpecialCharTok{$}\NormalTok{year[[}\DecValTok{2}\NormalTok{]], }\FunctionTok{aes}\NormalTok{(x,y, }\AttributeTok{color =} \StringTok{"97"}\NormalTok{)) }\SpecialCharTok{+}
  \FunctionTok{geom\_line}\NormalTok{(}\AttributeTok{data =}\NormalTok{ fit}\SpecialCharTok{$}\NormalTok{marginals.random}\SpecialCharTok{$}\NormalTok{year[[}\DecValTok{3}\NormalTok{]], }\FunctionTok{aes}\NormalTok{(x,y, }\AttributeTok{color =} \StringTok{"98"}\NormalTok{)) }
\end{Highlighting}
\end{Shaded}

\pandocbounded{\includegraphics[keepaspectratio]{GLMM_fit_files/figure-pdf/unnamed-chunk-7-2.pdf}}

\end{tcolorbox}

We now want to look at the hyperparameters of the model

\begin{tcolorbox}[enhanced jigsaw, opacityback=0, titlerule=0mm, toptitle=1mm, toprule=.15mm, opacitybacktitle=0.6, colframe=quarto-callout-warning-color-frame, arc=.35mm, rightrule=.15mm, title={Task}, colback=white, bottomrule=.15mm, bottomtitle=1mm, coltitle=black, left=2mm, leftrule=.75mm, breakable, colbacktitle=quarto-callout-warning-color!10!white]

Plot the posterior estimated of the standard deviation
\(\sigma_e = \sqrt{1/\tau_e}\) and
\(\sigma_\epsilon = \sqrt{1/\tau_\epsilon}\) of the two random effects.

Take hint

You can use the function \texttt{inla.tmarginal()} or the
\texttt{generate()} function

Click here to see the solution

\begin{Shaded}
\begin{Highlighting}[]
\CommentTok{\# Method 1 : use the inla.tmarginal() function}

\NormalTok{sd\_e }\OtherTok{=} \FunctionTok{inla.tmarginal}\NormalTok{(}\AttributeTok{fun =} \ControlFlowTok{function}\NormalTok{(x)}\FunctionTok{sqrt}\NormalTok{(}\DecValTok{1}\SpecialCharTok{/}\NormalTok{x), fit}\SpecialCharTok{$}\NormalTok{marginals.hyperpar}\SpecialCharTok{$}\StringTok{\textasciigrave{}}\AttributeTok{Precision for brood\_year}\StringTok{\textasciigrave{}}\NormalTok{)}

\NormalTok{sd\_eps }\OtherTok{=} \FunctionTok{inla.tmarginal}\NormalTok{(}\AttributeTok{fun =} \ControlFlowTok{function}\NormalTok{(x)}\FunctionTok{sqrt}\NormalTok{(}\DecValTok{1}\SpecialCharTok{/}\NormalTok{x), fit}\SpecialCharTok{$}\NormalTok{marginals.hyperpar}\SpecialCharTok{$}\StringTok{\textasciigrave{}}\AttributeTok{Precision for brood\_year\_chicken}\StringTok{\textasciigrave{}}\NormalTok{)}

\FunctionTok{ggplot}\NormalTok{() }\SpecialCharTok{+} \FunctionTok{geom\_line}\NormalTok{(}\AttributeTok{data =}\NormalTok{ sd\_e, }\FunctionTok{aes}\NormalTok{(x,y, }\AttributeTok{color =} \StringTok{"sd\_e"}\NormalTok{)) }\SpecialCharTok{+} \FunctionTok{geom\_line}\NormalTok{(}\AttributeTok{data =}\NormalTok{ sd\_eps, }\FunctionTok{aes}\NormalTok{(x,y, }\AttributeTok{color =} \StringTok{"sd\_epsilon"}\NormalTok{))}
\end{Highlighting}
\end{Shaded}

\pandocbounded{\includegraphics[keepaspectratio]{GLMM_fit_files/figure-pdf/unnamed-chunk-8-1.pdf}}

\begin{Shaded}
\begin{Highlighting}[]
\CommentTok{\# Method 1 : use the generate() function}


\CommentTok{\# we use this to get the correct name of the parameters}
\FunctionTok{bru\_names}\NormalTok{(fit)}
\end{Highlighting}
\end{Shaded}

\begin{verbatim}
                            height                               year 
                          "height"                             "year" 
                        brood_year                 brood_year_chicken 
                      "brood_year"               "brood_year_chicken" 
          Precision for brood_year   Precision for brood_year_chicken 
        "Precision_for_brood_year" "Precision_for_brood_year_chicken" 
\end{verbatim}

\begin{Shaded}
\begin{Highlighting}[]
\NormalTok{out }\OtherTok{=} \FunctionTok{generate}\NormalTok{(fit, }\AttributeTok{formula =} \SpecialCharTok{\textasciitilde{}} \FunctionTok{data.frame}\NormalTok{(}\AttributeTok{sd\_e =} \FunctionTok{sqrt}\NormalTok{(}\DecValTok{1}\SpecialCharTok{/}\NormalTok{Precision\_for\_brood\_year),}
                                     \AttributeTok{sd\_eps =} \FunctionTok{sqrt}\NormalTok{(}\DecValTok{1}\SpecialCharTok{/}\NormalTok{Precision\_for\_brood\_year\_chicken)),}
               \AttributeTok{n.samples =} \DecValTok{10000}\NormalTok{)}

\NormalTok{sd\_sample }\OtherTok{=} \FunctionTok{data.frame}\NormalTok{(}\AttributeTok{sd\_e =} \FunctionTok{unlist}\NormalTok{(}\FunctionTok{sapply}\NormalTok{(out, }\ControlFlowTok{function}\NormalTok{(x) x[}\DecValTok{1}\NormalTok{])),}
                       \AttributeTok{sd\_eps =} \FunctionTok{unlist}\NormalTok{(}\FunctionTok{sapply}\NormalTok{(out, }\ControlFlowTok{function}\NormalTok{(x) x[}\DecValTok{2}\NormalTok{])))}
                       

\FunctionTok{ggplot}\NormalTok{() }\SpecialCharTok{+} \FunctionTok{geom\_line}\NormalTok{(}\AttributeTok{data =}\NormalTok{ sd\_e, }\FunctionTok{aes}\NormalTok{(x,y, }\AttributeTok{color =} \StringTok{"sd\_e"}\NormalTok{)) }\SpecialCharTok{+} 
  \FunctionTok{geom\_density}\NormalTok{(}\AttributeTok{data =}\NormalTok{ sd\_sample, }\FunctionTok{aes}\NormalTok{(}\AttributeTok{x =}\NormalTok{ sd\_e, }\AttributeTok{color =} \StringTok{"sd\_e"}\NormalTok{), }\AttributeTok{linetype =} \StringTok{"dashed"}\NormalTok{) }\SpecialCharTok{+}
  \FunctionTok{geom\_line}\NormalTok{(}\AttributeTok{data =}\NormalTok{ sd\_eps, }\FunctionTok{aes}\NormalTok{(x,y, }\AttributeTok{color =} \StringTok{"sd\_epsilon"}\NormalTok{))}\SpecialCharTok{+} 
  \FunctionTok{geom\_density}\NormalTok{(}\AttributeTok{data =}\NormalTok{ sd\_sample, }\FunctionTok{aes}\NormalTok{(}\AttributeTok{x =}\NormalTok{ sd\_eps, }\AttributeTok{color =} \StringTok{"sd\_epsilon"}\NormalTok{), }\AttributeTok{linetype =} \StringTok{"dashed"}\NormalTok{)}
\end{Highlighting}
\end{Shaded}

\pandocbounded{\includegraphics[keepaspectratio]{GLMM_fit_files/figure-pdf/unnamed-chunk-8-2.pdf}}

\begin{Shaded}
\begin{Highlighting}[]
\CommentTok{\# Note that, if one is interested in hyperparameters, sampling is not always the best choice!}
\end{Highlighting}
\end{Shaded}

\end{tcolorbox}

\section{Change prior distributions}\label{change-prior-distributions}

Before trying to change the priors for the hyperparameters we can check
which priors are actually used in the model

\begin{Shaded}
\begin{Highlighting}[]
\FunctionTok{inla.priors.used}\NormalTok{(fit)}
\end{Highlighting}
\end{Shaded}

\begin{verbatim}
section=[family]
    tag=[INLA.Data1] component=[poisson]
section=[random]
    tag=[year] component=[year]
        group.theta1:
            parameter=[logit correlation]
            prior=[normal]
            param=[0.0, 0.2]
    tag=[] component=[]
    tag=[brood_year] component=[brood_year]
        theta1:
            parameter=[log precision]
            prior=[loggamma]
            param=[1e+00, 5e-05]
        group.theta1:
            parameter=[logit correlation]
            prior=[normal]
            param=[0.0, 0.2]
    tag=[brood_year_chicken] component=[brood_year_chicken]
        theta1:
            parameter=[log precision]
            prior=[loggamma]
            param=[1e+00, 5e-05]
        group.theta1:
            parameter=[logit correlation]
            prior=[normal]
            param=[0.0, 0.2]
section=[linear]
    tag=[height] component=[height]
        beta:
            parameter=[height]
            prior=[normal]
            param=[0.000, 0.001]
\end{verbatim}

From the output we see that the precision for the linear effect of
altitude is 0.001 (which means the sd is \(1/\sqrt{0.001} = 31.62\))

The precisions for the random effects have a Gamma prior with parameters
1 and 5e-05.

\subsection{Change the precision for the linear
effects}\label{change-the-precision-for-the-linear-effects}

The precision for linear effects is set in the component definition. For
example, if we want to increase the precision to 0.1 for we define the
relative components as:

\begin{Shaded}
\begin{Highlighting}[]
\NormalTok{cmp}\OtherTok{=} \ErrorTok{\textasciitilde{}} \SpecialCharTok{{-}}\DecValTok{1} \SpecialCharTok{+} \FunctionTok{year}\NormalTok{(YEAR, ... ) }\SpecialCharTok{+}
  \FunctionTok{height}\NormalTok{(cHEIGHT, }\AttributeTok{model =}  \StringTok{"linear"}\NormalTok{, }\AttributeTok{prec.linear =} \FloatTok{0.1}\NormalTok{) }\SpecialCharTok{+}\NormalTok{ ...}
\end{Highlighting}
\end{Shaded}

\begin{tcolorbox}[enhanced jigsaw, opacityback=0, titlerule=0mm, toptitle=1mm, toprule=.15mm, opacitybacktitle=0.6, colframe=quarto-callout-warning-color-frame, arc=.35mm, rightrule=.15mm, title={Task}, colback=white, bottomrule=.15mm, bottomtitle=1mm, coltitle=black, left=2mm, leftrule=.75mm, breakable, colbacktitle=quarto-callout-warning-color!10!white]

Run the model again using 0.1 as default precision for the altitude
slope parameter and for the \texttt{year} fixed effect

Click here to see the solution

\begin{Shaded}
\begin{Highlighting}[]
\NormalTok{cmp2 }\OtherTok{=}  \ErrorTok{\textasciitilde{}} \SpecialCharTok{{-}}\DecValTok{1} \SpecialCharTok{+} \FunctionTok{year}\NormalTok{(YEAR, }\AttributeTok{model =} \StringTok{"iid"}\NormalTok{, }\AttributeTok{initial =} \FunctionTok{log}\NormalTok{(}\FloatTok{0.1}\NormalTok{), }\AttributeTok{fixed =}\NormalTok{ T) }\SpecialCharTok{+}
  \FunctionTok{height}\NormalTok{(cHEIGHT, }\AttributeTok{model =}  \StringTok{"linear"}\NormalTok{, }\AttributeTok{prec.linear =} \FloatTok{0.1}\NormalTok{) }\SpecialCharTok{+} 
  \FunctionTok{brood\_year}\NormalTok{(ij , }\AttributeTok{model =} \StringTok{"iid"}\NormalTok{) }\SpecialCharTok{+}
  \FunctionTok{brood\_year\_chicken}\NormalTok{(ijk, }\AttributeTok{model=} \StringTok{"iid"}\NormalTok{)}

\NormalTok{fit2 }\OtherTok{=} \FunctionTok{bru}\NormalTok{(cmp2, lik) }
\end{Highlighting}
\end{Shaded}

Note that we can use the same observation model as before since both the
formula and the dataset are unchanged.

\end{tcolorbox}

\textbf{Changing the precision for the random effects}

Priors on the hyperparameters of the random effects model must be passed
by defining argument \texttt{hyper} within component of interest

\begin{Shaded}
\begin{Highlighting}[]
\CommentTok{\# First we define the logGamma (0.01,0.01) prior }

\NormalTok{prec.prior }\OtherTok{\textless{}{-}} \FunctionTok{list}\NormalTok{(}\AttributeTok{prec =} \FunctionTok{list}\NormalTok{(}\AttributeTok{prior =} \StringTok{"loggamma"}\NormalTok{,  }\CommentTok{\# prior name}
                               \AttributeTok{param =} \FunctionTok{c}\NormalTok{(}\FloatTok{0.01}\NormalTok{, }\FloatTok{0.01}\NormalTok{))) }\CommentTok{\# prior parameters}
                        

\NormalTok{cmp3 }\OtherTok{=}  \ErrorTok{\textasciitilde{}} \SpecialCharTok{{-}}\DecValTok{1} \SpecialCharTok{+} \FunctionTok{year}\NormalTok{(YEAR, }\AttributeTok{model =} \StringTok{"iid"}\NormalTok{, }\AttributeTok{initial =} \FunctionTok{log}\NormalTok{(}\FloatTok{0.1}\NormalTok{), }\AttributeTok{fixed =}\NormalTok{ T) }\SpecialCharTok{+}
  \FunctionTok{height}\NormalTok{(cHEIGHT, }\AttributeTok{model =}  \StringTok{"linear"}\NormalTok{, }\AttributeTok{prec.linear =} \FloatTok{0.1}\NormalTok{) }\SpecialCharTok{+} 
  \FunctionTok{brood\_year}\NormalTok{(ij , }\AttributeTok{model =} \StringTok{"iid"}\NormalTok{, }\AttributeTok{hyper =}\NormalTok{ prec.prior) }\SpecialCharTok{+}
  \FunctionTok{brood\_year\_chicken}\NormalTok{(ijk, }\AttributeTok{model=} \StringTok{"iid"}\NormalTok{, }\AttributeTok{hyper =}\NormalTok{ prec.prior)}


\NormalTok{fit3 }\OtherTok{=} \FunctionTok{bru}\NormalTok{(cmp3, lik) }
\end{Highlighting}
\end{Shaded}

Note that for \texttt{model\ =\ "linear"} the value for the precision is
expressed in natural scale while for
\texttt{model\ =\ "iid\textquotesingle{}} (and all other random effects)
the value is expressed in log scale.

\subsubsection{Visualizing the posterior
marginals}\label{visualizing-the-posterior-marginals}

Posterior marginal distributions of the fixed effects parameters and the
hyperparameters can be visualized also using the \texttt{plot()}
function by calling the name of the component. For example, if want to
visualize the posterior density of the slope \(\beta\) we can type:

\begin{Shaded}
\begin{Highlighting}[]
\FunctionTok{plot}\NormalTok{(fit, }\StringTok{"height"}\NormalTok{)}
\end{Highlighting}
\end{Shaded}

\begin{center}
\pandocbounded{\includegraphics[keepaspectratio]{GLMM_fit_files/figure-pdf/unnamed-chunk-13-1.pdf}}
\end{center}

\begin{tcolorbox}[enhanced jigsaw, opacityback=0, titlerule=0mm, toptitle=1mm, toprule=.15mm, opacitybacktitle=0.6, colframe=quarto-callout-warning-color-frame, arc=.35mm, rightrule=.15mm, title={Task}, colback=white, bottomrule=.15mm, bottomtitle=1mm, coltitle=black, left=2mm, leftrule=.75mm, breakable, colbacktitle=quarto-callout-warning-color!10!white]

Plot the posterior marginals for the precision of brood and individual
random effect \(\pi(\tau_e|y)\)

Take hint

See the \texttt{summary()} output to check the names for the different
model components.

Click here to see the solution

\begin{Shaded}
\begin{Highlighting}[]
\FunctionTok{plot}\NormalTok{(fit, }\StringTok{"Precision for brood\_year"}\NormalTok{)}
\end{Highlighting}
\end{Shaded}

\pandocbounded{\includegraphics[keepaspectratio]{GLMM_fit_files/figure-pdf/unnamed-chunk-14-1.pdf}}

\end{tcolorbox}




\end{document}
